\section{Przygotowanie i uzasadnienie potrzeby realizacji}
\label{sec:analiza}

\subsection{Grupa projektowa i podział obowiązków}

Projekt realizuje trzyosobowa grupa studentów kierunku Informatyka techniczna na Wydziale Informatyki i~Telekomunikacji Politechniki Wrocławskiej. Prowadzącym przedmiot jest dr~inż. Roman Ptak. Skład grupy zestawiono w~tabeli~\ref{tab:grupa}. Podział obowiązków zawiera tabela~\ref{tab:obowiazki}. Na początku są obowiązki trwające przez cały projekt, dalej bloki etapów. Przy każdym obowiązku wskazano jedną osobę odpowiedzialną.

\begin{table}[H]
	\centering
	\caption{Skład grupy projektowej}
	\label{tab:grupa}
	\begin{tabular}{@{}ll@{}}
		\toprule
		Imię i~nazwisko & Numer indeksu \\
		\midrule
		Karol Gąsior & 280913 \\
		Mikołaj Winiecki & 280952 \\
		Mateusz Złośnik & 281201 \\
		\bottomrule
	\end{tabular}
\end{table}

\begin{longtable}{@{}p{0.72\textwidth}l@{}}
	\caption{Podział obowiązków w~grupie projektowej}
	\label{tab:obowiazki}\\
	\toprule
	Obowiązek & Odpowiedzialny \\
	\midrule
	\endfirsthead
	\caption[]{Podział obowiązków w~grupie projektowej (cd.)}\\
	\toprule
	Obowiązek & Odpowiedzialny \\
	\midrule
	\endhead
	\bottomrule
	\endfoot
	\multicolumn{2}{@{}l}{\textbf{Obowiązki ogólne}} \\
	\midrule
	Prowadzenie repozytorium i~opis zmian & Karol Gąsior \\
	Kompilacja dokumentacji i~kontrola składu PDF & Mikołaj Winiecki \\
	Pilnowanie terminów oraz kontaktu z~prowadzącym & Mateusz Złośnik \\
	Czytanie rozdziałów przed oddaniem etapu & Karol Gąsior \\
	Przegląd kodu przed uznaniem funkcji za gotową & Mikołaj Winiecki \\
	Jedna konwencja nazw w~modelach, widokach i~w~tekście & Mateusz Złośnik \\
	Kopie pliku bazy oraz danych testowych & Karol Gąsior \\
	Wpis ustaleń grupy do dokumentacji & Mikołaj Winiecki \\
	Pilnowanie zakresu MVP: jeden sklep, bez płatności i~bez aplikacji mobilnej & Mateusz Złośnik \\
	\midrule
	\multicolumn{2}{@{}l}{\textbf{Etap 1. Analiza}} \\
	\midrule
	Opis słowny sklepu: asortyment, dostawy, sprzedaż i~stany & Karol Gąsior \\
	Problemy prowadzenia sklepu bez systemu & Mateusz Złośnik \\
	Interesariusze i~otoczenie podsystemu & Mikołaj Winiecki \\
	Analiza SWOT & Karol Gąsior \\
	Czynniki krytyczne sukcesu & Mateusz Złośnik \\
	Cel systemu i~granice zakresu MVP & Mikołaj Winiecki \\
	Wykonalność techniczna i~szacunek pracochłonności & Karol Gąsior \\
	Zgodność harmonogramu z~terminami laboratorium & Mikołaj Winiecki \\
	\midrule
	\multicolumn{2}{@{}l}{\textbf{Etap 2. Projekt}} \\
	\midrule
	Wymagania funkcjonalne MVP & Mikołaj Winiecki \\
	Wymagania niefunkcjonalne & Mateusz Złośnik \\
	Architektura trójwarstwowa: przeglądarka, Django, SQLite & Karol Gąsior \\
	Diagram przypadków użycia & Mateusz Złośnik \\
	Diagram klas albo komponentów & Mikołaj Winiecki \\
	Diagram związków encji (ERD) & Karol Gąsior \\
	Klucze, związki i~uzasadnienie postaci normalnej & Mikołaj Winiecki \\
	Słownik danych: atrybuty, typy i~ograniczenia & Mateusz Złośnik \\
	\midrule
	\multicolumn{2}{@{}l}{\textbf{Etap 3. Implementacja}} \\
	\midrule
	Szkielet projektu Django i~aplikacji sklepu & Mateusz Złośnik \\
	Model produktu i~stanu magazynowego & Karol Gąsior \\
	Model dostawy i~pozycji dostawy & Mikołaj Winiecki \\
	Model sprzedaży i~pozycji sprzedaży & Karol Gąsior \\
	Migracje schematu do SQLite & Mikołaj Winiecki \\
	Widoki i~formularze produktów & Mateusz Złośnik \\
	Widoki i~formularze dostaw oraz sprzedaży, ze zmianą stanu & Karol Gąsior \\
	Konfiguracja panelu administratora & Mikołaj Winiecki \\
	Test przejścia CRUD i~opis wyników & Mateusz Złośnik \\
	\midrule
	\multicolumn{2}{@{}l}{\textbf{Etap 4. Wdrożenie}} \\
	\midrule
	Plik zależności i~środowisko wirtualne & Mateusz Złośnik \\
	Instrukcja uruchomienia na komputerze lokalnym & Karol Gąsior \\
	Dane początkowe asortymentu i~stanów & Mikołaj Winiecki \\
	Konta do pracy codziennej i~do pokazu & Mateusz Złośnik \\
	Kopia zapasowa bazy SQLite i~jej odtworzenie & Karol Gąsior \\
	Checklista odbioru po wdrożeniu & Mikołaj Winiecki \\
	\midrule
	\multicolumn{2}{@{}l}{\textbf{Etap 5. Prezentacja i~zaliczenie}} \\
	\midrule
	Scenariusz prezentacji i~podział wypowiedzi & Mateusz Złośnik \\
	Dane pokazowe w~działającym sklepie & Karol Gąsior \\
	Demonstracja funkcji MVP & Mikołaj Winiecki \\
	Konspekt albo slajdy & Mateusz Złośnik \\
	Próba generalna i~pomiar czasu & Karol Gąsior \\
	Wnioski w~dokumentacji i~oddanie PDF & Mikołaj Winiecki \\
\end{longtable}

\subsection{Opis słowny sklepu}
\label{sec:opis-sklepu}

Przedmiotem opisu jest jeden lokalny sklep warzywny. Sklep sprzedaje warzywa i~owoce klientom detalicznym, przy ladzie. Nie prowadzi sprzedaży wysyłkowej, nie obsługuje wielu punktów i~nie wydaje kart lojalnościowych. Płatność odbywa się poza systemem informatycznym.

Asortyment składa się z~produktów spożywczych sprzedawanych na wagę albo na sztuki. Każdy produkt ma nazwę, jednostkę miary oraz aktualną cenę sprzedaży. Lista nie jest zamknięta: w~ciągu roku sklep dopisuje pozycję sezonową, zmienia cenę i~wycofuje towar, którego już nie sprowadza.

Dostawy przychodzą od dostawców. Jedna dostawa dotyczy jednego dostawcy, ma datę przyjęcia i~może zawierać wiele produktów. Pozycja dostawy wskazuje produkt oraz przyjętą ilość. Przyjęcie towaru zwiększa stan magazynowy tego produktu.

Sprzedaż jest rejestrowana jako jedno zdarzenie z~jedną lub wieloma pozycjami. Pozycja sprzedaży wskazuje produkt oraz sprzedaną ilość, liczoną według ceny obowiązującej w~chwili sprzedaży. Sprzedaż zmniejsza stan magazynowy. Nie sprzedaje się ilości większej niż aktualny stan.

Stan magazynowy to aktualna ilość danego produktu w~tym jednym sklepie. Stan rośnie wyłącznie przez dostawę i~maleje wyłącznie przez sprzedaż. Brak towaru oznacza stan równy zero do czasu kolejnej dostawy.

\subsection{Problemy prowadzenia sklepu bez systemu}
\label{sec:problemy}
% TODO: obowiązek 1.2

\subsection{Interesariusze i otoczenie podsystemu}
\label{sec:interesariusze}

Podsystem obejmuje ewidencję asortymentu, dostaw, sprzedaży i~stanów w~jednym sklepie. Korzystają z~niego osoby pracujące na miejscu. Dostawca i~klient zostają w~otoczeniu: nie mają konta i~nie wprowadzają danych. W~małym sklepie osoba prowadząca i~sprzedawca mogą być tą samą osobą. W~projekcie rozdziela się role, a~nie etaty. Zestawienie zawiera tabela~\ref{tab:interesariusze}.

\begin{table}[H]
	\centering
	\caption{Interesariusze podsystemu sklepu}
	\label{tab:interesariusze}
	\begin{tabular}{@{}L{0.28\textwidth}L{0.22\textwidth}L{0.40\textwidth}@{}}
		\toprule
		Interesariusz & Relacja & Oczekiwanie \\
		\midrule
		Osoba prowadząca sklep & użytkownik & aktualny asortyment, ceny i~stany \\
		Sprzedawca & użytkownik & szybki zapis dostawy i~sprzedaży \\
		Dostawca & otoczenie & brak konta; dostawa jest tylko rejestrowana \\
		Klient detaliczny & otoczenie & obsługa przy ladzie, bez konta i~bez płatności w~systemie \\
		\bottomrule
	\end{tabular}
\end{table}

Otoczenie stanowi jeden lokal, sezonowa zmienność asortymentu oraz płatność rozliczana poza aplikacją. Podsystem nie łączy się z~kasą, ze sklepem internetowym ani z~innym programem magazynowym. Granicą jest zapis przyjęcia towaru i~sprzedaży przy ladzie.

\subsection{Analiza SWOT}
\label{sec:swot}

Analizę SWOT odniesiono do planowanego systemu, a~nie do rynku warzyw. Wynik zestawia tabela~\ref{tab:swot}.

\begin{table}[H]
	\centering
	\caption{Analiza SWOT planowanego systemu}
	\label{tab:swot}
	\begin{tabular}{@{}L{0.22\textwidth}L{0.70\textwidth}@{}}
		\toprule
		Czynnik & Ustalenie \\
		\midrule
		Mocne strony &
		Jeden sklep i~cztery operacje: produkt, stan, dostawa, sprzedaż. Reguła stanu jest jednoznaczna. Django daje formularze, panel administratora i~migracje, a~SQLite nie wymaga osobnego serwera bazy. \\
		\addlinespace
		Słabe strony &
		Na etapie analizy nie ma jeszcze kodu. SQLite słabo znosi wielu jednoczesnych zapisów, co przy jednym stanowisku jest do przyjęcia. System nie liczy płatności. \\
		\addlinespace
		Szanse &
		Stany wynikają z~zapisanych ruchów, a~nie z~pamięci albo kartek. Sezonowa zmiana asortymentu sprowadza się do dopisania produktu i~ceny. Ten sam zakres wystarcza do codziennej pracy i~do zaliczenia. \\
		\addlinespace
		Zagrożenia &
		Implementacja i~wdrożenie przypadają na jeden termin, 30.11.2026. Dołożenie płatności, lojalności albo drugiego sklepu zabiera czas przeznaczony na stany. Błąd przy zmianie ilości psuje wiarygodność ewidencji. \\
		\bottomrule
	\end{tabular}
\end{table}

O utrzymaniu się w~terminie decydują wąski zakres oraz poprawna zmiana stanu przy dostawie i~sprzedaży.

\subsection{Czynniki krytyczne sukcesu}
\label{sec:csf}

Czynnik krytyczny sukcesu to warunek, bez którego system nie spełnia celu z~podrozdziału~\ref{sec:cel}. Lista jest krótka i~da się ją sprawdzić na działającej aplikacji. Zestawienie zawiera tabela~\ref{tab:csf}.

\begin{table}[H]
	\centering
	\caption{Czynniki krytyczne sukcesu}
	\label{tab:csf}
	\begin{tabular}{@{}L{0.48\textwidth}L{0.44\textwidth}@{}}
		\toprule
		Czynnik & Skutek braku \\
		\midrule
		Stan wynika tylko z~dostawy i~sprzedaży i~nie schodzi poniżej zera &
		ewidencja rozmija się z~półką \\
		\addlinespace
		Dostawę i~sprzedaż da się zapisać przy ladzie, w~kilku polach &
		obsługa wraca do notatek \\
		\addlinespace
		Zakres zostaje przy jednym sklepie, bez płatności, lojalności i~aplikacji mobilnej &
		prace nie mieszczą się w~terminie 30.11.2026 \\
		\addlinespace
		Na 14.12.2026 działa lokalnie ścieżka: produkt, dostawa, sprzedaż, widoczny stan &
		brak materiału na zaliczenie \\
		\addlinespace
		Plik bazy da się skopiować i~odtworzyć &
		pomyłka przy wpisie kasuje dane pokazu \\
		\bottomrule
	\end{tabular}
\end{table}

\subsection{Cel systemu i zakres}
\label{sec:cel}

Celem systemu jest prowadzenie asortymentu, dostaw, sprzedaży i~stanów jednego lokalnego sklepu warzywnego. Aktualna ilość produktu ma wynikać z~zapisanych ruchów. Osoba przy ladzie ma widzieć, czy towar jest, i~mieć możliwość dopisania dostawy albo sprzedaży bez osobnej kartki.

W~zakresie zostają dane jednego sklepu: produkt z~nazwą, jednostką i~ceną, stan magazynowy, dostawa z~pozycjami oraz sprzedaż z~pozycjami. Obsługa idzie przez przeglądarkę na komputerze lokalnym, formularzami albo panelem administratora Django. Baza to jeden plik SQLite.

Poza zakresem zostają: wiele punktów sprzedaży, płatności i~kasa, program lojalnościowy, aplikacja mobilna, sprzedaż wysyłkowa oraz składanie zamówień u~dostawcy. Dostawa w~systemie oznacza przyjęcie towaru, który już dotarł. Granicę zakresu zestawia tabela~\ref{tab:zakres}.

\begin{table}[H]
	\centering
	\caption{Granica zakresu MVP}
	\label{tab:zakres}
	\begin{tabular}{@{}L{0.46\textwidth}L{0.46\textwidth}@{}}
		\toprule
		W zakresie & Poza zakresem \\
		\midrule
		Produkt, cena, jednostka & wiele sklepów \\
		Stan magazynowy & płatności i kasa \\
		Dostawa i pozycje dostawy & program lojalnościowy \\
		Sprzedaż i pozycje sprzedaży & aplikacja mobilna \\
		Formularz albo panel administratora & sprzedaż wysyłkowa \\
		Jedna baza SQLite na jednym komputerze & zamówienia do dostawcy \\
		\bottomrule
	\end{tabular}
\end{table}

\subsection{Wykonalność}
\label{sec:wykonalnosc}

Projekt jest wykonalny technicznie na ustalonym stosie. Django i~SQLite uruchamiają się lokalnie, bez osobnego serwera bazy i~bez usług zewnętrznych. Schemat powstaje migracjami, a~dostęp do danych idzie przez ORM. Panel administratora pokrywa wprowadzanie produktów, dostaw i~sprzedaży także wtedy, gdy własne formularze powstają później. Reguła biznesowa jest jedna: dostawa zwiększa stan, sprzedaż go zmniejsza i~nie przekracza stanu. Przy jednym stanowisku ograniczenie jednoczesnego zapisu w~SQLite nie blokuje pracy.

Szacunek pracochłonności jest względny. Nie mierzono godzin. Wynika z~kalendarza laboratorium oraz z~tego, że implementacja modeli, formularzy i~zmiany stanu jest jedynym blokiem zależnym od wcześniejszego schematu. Porównanie zawiera tabela~\ref{tab:pracochlonnosc}. Trzy osoby mogą równolegle zbierać wymagania i~rysować diagramy. Widoki powstają po modelach. Wdrożenie nie dokłada serwera: zostaje środowisko na komputerze, plik bazy i~instrukcja uruchomienia.

\begin{table}[H]
	\centering
	\caption{Względny szacunek pracochłonności}
	\label{tab:pracochlonnosc}
	\begin{tabular}{@{}L{0.24\textwidth}L{0.22\textwidth}L{0.44\textwidth}@{}}
		\toprule
		Blok & Udział & Na czym polega \\
		\midrule
		Analiza & mniejszy & opis sklepu, cel i zakres \\
		Projekt & średni & wymagania, architektura, ERD \\
		Implementacja & największy & modele, stan, formularze \\
		Wdrożenie & mniejszy & start lokalny i kopia bazy \\
		Prezentacja & mniejszy & pokaz gotowej ścieżki MVP \\
		\bottomrule
	\end{tabular}
\end{table}

Najciaśniejszy termin to 30.11.2026, bo obejmuje implementację i~wdrożenie. Jeśli własne formularze się opóźnią, panel administratora nadal pozwala pokazać zapis produktu, dostawy, sprzedaży i~wynikowy stan.

\subsection{Plan projektu i harmonogram}

\begin{table}[H]
	\centering
	\caption{Harmonogram etapów projektu}
	\label{tab:harmonogram}
	\begin{tabular}{@{}lll@{}}
		\toprule
		Data & Etap & Rezultat \\
		\midrule
		05.10.2026 & 0 & Wprowadzenie do projektu \\
		19.10.2026 & 1 & Analiza (start) \\
		17.11.2026 & 2 & Faza projektowa \\
		30.11.2026 & 3+4 & Implementacja i wdrożenie \\
		14.12.2026 & 5 & Prezentacja i zaliczenie \\
		\bottomrule
	\end{tabular}
\end{table}
